% !TEX program = pdflatex
\documentclass[10pt]{article}
% Shared layout for the Spanish and English CVs (build with cv/build.sh).
\usepackage[utf8]{inputenc}
\usepackage[T1]{fontenc}
\usepackage{lmodern}
% Clean copy/paste and ATS text extraction
\input{glyphtounicode}
\pdfgentounicode=1
\usepackage[margin=0.55in]{geometry}
\usepackage{enumitem}
\usepackage{titlesec}
\usepackage{parskip}
\setlength{\parskip}{0.2em}
\usepackage[hidelinks]{hyperref}

\titleformat{\section}{\large\bfseries}{}{0pt}{}[\vspace{-0.6em}\rule{\linewidth}{0.4pt}]
\titlespacing{\section}{0pt}{0.4em}{0.2em}
\setlist[itemize]{leftmargin=1.2em, itemsep=0.05em, topsep=0.1em, parsep=0pt}
\pagestyle{empty}

\newcommand{\job}[4]{\noindent\textbf{#1} \textbar{} #2 \textperiodcentered{} \textit{#4} \hfill #3}
\newcommand{\project}[3]{\noindent\textbf{\href{#2}{#1}} \textbar{} \textit{#3}\\}

\hypersetup{
  pdftitle={Miguel Ángel Altamar Rodríguez - CV},
  pdfauthor={Miguel Ángel Altamar Rodríguez},
  pdfsubject={Desarrollador Full-Stack | Aplicaciones Web, APIs e Integraciones},
  pdfkeywords={Full-Stack, React, Next.js, Angular, TypeScript, Node.js, NestJS, Fastify, Python, FastAPI, Django, Java, Spring Boot, C\#, ASP.NET, API REST, GraphQL, SOAP, WS-Security, Microservicios, PostgreSQL, Docker, CI/CD, AWS, Azure, Business Central, AL}
}

\begin{document}

\begin{center}
  {\LARGE\bfseries Miguel Ángel Altamar Rodríguez}\\[0.3em]
  {\large Desarrollador Full-Stack | Aplicaciones Web, APIs e Integraciones}\\[0.4em]
  Barranquilla, Colombia \enspace\textbullet\enspace
  \href{mailto:maltamarr@outlook.com}{maltamarr@outlook.com} \enspace\textbullet\enspace
  +57 300 512 3022\\[0.15em]
  \href{https://migueiangel.github.io/portfolio/es/}{Portafolio} \enspace\textbullet\enspace
  \href{https://www.linkedin.com/in/mangel2002/}{LinkedIn} \enspace\textbullet\enspace
  \href{https://github.com/MigueIAngel}{GitHub}
\end{center}
\vspace{0.1em}

\section*{Perfil profesional}
Ingeniero de sistemas full-stack que disfruta llevar una funcionalidad de punta a punta. Construyo aplicaciones web con React, Angular y Next.js sobre APIs en Node.js, Python, Java y .NET, y extiendo e integro Microsoft Dynamics 365 Business Central: facturación electrónica con proveedores de certificación fiscal (PAC), APIs de integración y reportes.

\section*{Experiencia profesional}
\job{Desarrollador AL Senior}{4 Ways Tech}{Mar 2026 – Presente}{Remoto}
\begin{itemize}
  \item \textbf{Facturación electrónica e-CF (República Dominicana)} con el PAC SERES desde Business Central: XML fiscal de Facturas de Crédito Fiscal (31), de Consumo (32) y notas de crédito, y servicio SOAP con WS-Security para envío y consulta de estado.
  \item \textbf{14 reportes RDL/RDLC} construidos desde cero para 8 clientes (arqueo de caja, antigüedad de CxC/CxP, balance de comprobación, certificados de retención, entre otros) y \textbf{13 formatos} de facturas, órdenes de compra y cheques mejorados.
  \item \textbf{APIs de integración}: páginas API para sistemas externos (artículos, unidades de medida, valores de dimensión) con una tabla y una codeunit de log que dan trazabilidad a cada llamada.
  \item \textbf{Soporte en producción}: corrección de defectos de extensiones AL (dimensiones, permisos, archivos bancarios) y refactorización de código obsoleto tras actualizaciones.
\end{itemize}
\job{Desarrollador Business Central}{LLB Solutions}{Jun 2025 – Mar 2026}{Barranquilla, Colombia}
\begin{itemize}
  \item Diseño, desarrollo e implementación de soluciones a la medida en Business Central según las necesidades de cada cliente, junto a consultores funcionales.
  \item Extensión del ERP en AL con reportes, páginas e integraciones con sistemas externos, y pruebas automatizadas para reducir regresiones.
\end{itemize}
\job{Desarrollador Fullstack}{Wizybot}{Abr 2025 – May 2025}{Barranquilla, Colombia}
\begin{itemize}
  \item Diseño e integración de APIs REST entre herramientas internas y servicios de terceros.
\end{itemize}

\section*{Proyectos destacados}
\begin{itemize}
  \item \textbf{\href{https://github.com/MigueIAngel/orderflow}{OrderFlow}} \textit{(NestJS, FastAPI, Fastify, Redis Streams, React, Docker)}: cinco microservicios con saga coreografiado, outbox transaccional, consumidores idempotentes, compensación y API gateway con circuit breakers; prueba end-to-end del saga en CI.
  \item \textbf{\href{https://github.com/MigueIAngel/finance-tracker}{Finance Tracker}} \textit{(Fastify, Drizzle ORM, PostgreSQL, Zod, Next.js)}: API REST protegida por API key detrás de rutas proxy en el servidor; clasificación de transacciones con un modelo de lenguaje.
  \item \textbf{\href{https://github.com/MigueIAngel/docuchat-ai}{DocuChat AI}} \textit{(FastAPI, ChromaDB, Gemini, React)}: RAG sobre PDFs con respuestas en streaming (SSE) y citas por página.
  \item \textbf{\href{https://github.com/MigueIAngel/linkvault}{LinkVault}} \textit{(Next.js 16, Server Actions, Prisma, Auth.js)}: actualizaciones optimistas, OAuth de GitHub, i18n y protección SSRF.
\end{itemize}

\section*{Habilidades técnicas}
\begin{itemize}
  \item \textbf{Front-end:} React, Next.js, Angular, TypeScript, Tailwind CSS, Astro, HTMX, TanStack Query
  \item \textbf{Back-end:} Node.js, NestJS, Fastify, FastAPI, Django, Spring Boot, JEE/JPA, ASP.NET
  \item \textbf{Lenguajes:} TypeScript, JavaScript, Python, Java, C\#, SQL, AL
  \item \textbf{Integraciones y datos:} REST, GraphQL, SOAP, WS-Security, JWT, OAuth 2.0, XML, RAG; PostgreSQL, SQL Server, MySQL, DB2, DynamoDB, Redis
  \item \textbf{Arquitectura y DevOps:} microservicios, eventos, patrones de diseño, pruebas unitarias y de integración, Docker, GitHub Actions, CI/CD, AWS, Azure
\end{itemize}

\section*{Educación y certificaciones}
\noindent\textbf{Ingeniería de Sistemas} \textbar{} Universidad del Norte \textemdash{} \textit{Estudiante distinguido y becario} \hfill 2024\\
\href{https://learn.microsoft.com/api/credentials/share/es-es/MiguelAltamar-5421/468B2D2A0F8511E3?sharingId=88DCDE63C9CD6D93}{Microsoft Certified: Dynamics 365 Business Central Developer Associate (MB-820)} \textemdash{} Microsoft, 2025\\
\textbf{Idiomas:} Español (nativo), Inglés (B2)

\end{document}
